\documentclass{article}
\usepackage[utf8]{inputenc}
\usepackage[T1]{fontenc}
\usepackage{geometry}
\geometry{
    a4paper,
    left=2.5cm,
    right=2.5cm,
    top=2.5cm,
    bottom=2.5cm
}
\usepackage{graphicx}
\usepackage{float}
\usepackage{hyperref}
\usepackage[english]{babel}
\usepackage{amsmath}
\usepackage{amssymb}
\usepackage{array}
\usepackage{booktabs}
\usepackage{caption}
\usepackage{subcaption}
\usepackage{enumitem}
\usepackage{fancyhdr}
\usepackage{titlesec}
\usepackage{url}
\usepackage{textcomp}
\usepackage{xcolor}
\usepackage{listings}
\usepackage{multirow}
\usepackage{longtable}
\usepackage{tabularx}
\usepackage{times}

% Set up page style
\pagestyle{fancy}
\fancyhf{}
\rfoot{\thepage}
\renewcommand{\headrulewidth}{0pt}

% Title formatting
\titleformat{\chapter}[display]{\normalfont\bfseries}{\chaptertitle\thechapter}{0pt}{\Huge}
\titlespacing*{\chapter}{0pt}{50pt}{40pt}

% Document info
\title{Template Document}
\author{Academic Support Agent}
\date{\today}

\begin{document}
\maketitle

\title{Formatos en \LaTeX}
\author{Adolfo Ant n Bravo}
\maketitle{}
\section{El formato de fuente}
\label{sec:formato}

\subsection{Caracteres especiales}
\label{sec:caracteres}
\LaTeX distingue entre may sculas y min sculas cuando escribimos intrucciones, as  que hay que tener cuidado al escribirlas. Por ejemplo, \textbackslash LaTeX sirve para insertar el logotipo de \LaTeX en nuestros documentos.
Tambi n podemos escribir ``entre comillas'' -aunque al principio resulta un poco extra o-. Existen adema s muchos otros s mbolos que podemos a adir: \%, \$, \copyright, \ldots

\subsection{Espacios y saltos de l nea}
\label{sec:espacioslinea}

Tambi n es posible a adir espacios \hspace{2cm} entre palabras, aunque no suele utilizarse demasiado.
Para saltar de p rrafo utilizamos la orden \texttt{\textbackslash par} o simplemente dejamos una l nea en blanco.

\subsection{El tama o de la fuente}
\label{sec:fuente}

Tambi n podemos utilizar distintos tama os de fuente que siempre depender n de la clase del documento y de su configuraci n.

Existe letra \tiny min scula, \scriptsize tama o nota \footnotesize pie de p gina, \small peque a, \normalsize normal, \large grande, \Large m s grande, \LARGE todav a m s grande, \huge enorme y \Huge s per enorme. 

\subsection{El estilo de la fuente}
\label{sec:estilo}
\normalsize
Adem s, podemos escribir en \textbf{negrita}, \textit{cursiva}, \emph{enfatizado}, \uppercase{todo en may scula}, \textsc{capitalizado o versales}, \texttt{tama o fijo}, \ldots

\subsection{Sub ndices y super ndices}
\label{sec:subysuper}

La f rmula del  cido sulf rico es H\textsubscript{2}SO\textsubscript{4}. Y la famosa f rmula de Einstein es \textit{E=mc\textsuperscript{2}}

\subsection{Color}
\label{sec:color}
Para escribir con color debemos utilizar el paquete \color{red}\huge color
\normalsize

\color{black}
\section{El formato de p rrafo}
\label{sec:formatoparrafo}

\subsection{Cajas de texto}
\label{sec:cajastexto}
\mbox{Mi texto normal}

\makebox[10cm][r]{Mi texto alineado a la derecha}

\makebox[10cm][s]{Mi texto estirado para ocupar la caja}

\framebox[5cm][c]{Texto centrado con borde}

\framebox[10cm][s]{Una caja que ocupa exactamente lo que ocupa el texto}

\framebox[1.5\width][c]{Una caja que ocupa 1,5 veces m s que lo que ocupa el texto}

\framebox[0.5\width][l]{Una caja que ocupa la mitad (0.5) de lo que ocupa el texto}

\framebox[0.3\linewidth][c]{Una caja que ocupa un tercio (0.3) de la p gina}

\subsection{Bordes y sombreado de cajas}
\label{sec:bordesysombreadocajas}
{\fboxsep 8pt \fboxrule 1pt \fbox{
    \parbox[b]{10cm}{Esta es una caja con borde de 1 punto y margen de 8 puntos}
    } % Cierra fbox
} % cierra el grupo

\subsection{Color de las cajas}
\label{sec:colorcajas}

Podemos colocar las cajas con \textbf{fcolorbox}:

\colorbox{red}{texto escrito en negro dentro de una caja de color rojo}
\\

\fcolorbox{blue}{white}{
    \parbox[b]{.2\linewidth}{texto escrito en negro dentro de una caja de p rrafo de color negro con borde azul}
}


\subsection{Alineaci n de las cajas}
\label{sec:alineacioncajas}

Para mostrar la alineaci n de las cajas utilizaremos cajas de p rrafo (\texttt{parbox} con ancho de \texttt{0.3 linewidth}:

\parbox[b]{0.3\linewidth}{\raggedright{Caja de p rrafo a la izqda. con texto alineado a la izqda.}}
\parbox[b]{0.3\linewidth}{\centering Caja de p rrafo centrada con texto centrado}
\parbox[b]{0.3\linewidth}{\raggedleft{Caja de p rrafo a la dcha. con texto alineado a la dcha.}}

\subsection{Espaciado entre p rrafos}
\label{sec:espaciadoparrafos}

Podemos saltar de p rrafo dejando una l nea en blanco con \texttt{par}\par Y de p gina con con \texttt{newpage}\newpage{}
 sta es otra p gina. y tambi n podemos saltar de l nea sin saltar de p rrafo con \texttt{newline} o dos barras invertidas.\newline{} Esto es un salto de l nea. Puede observarse que en los saltos de l nea no se indenta mientras que en los saltos de p rrafo, s .

Entre los p rrafos puede haber espacios peque os:\smallskip{}

Espacios medios:
\medskip{}

O espacios grandes:
\bigskip{}

Tambi n podemos dejar el espacio que queramos con \texttt{vspace}\vspace{4cm}

Acabamos de dejar 4 cm. de espacio.
\subsection{Letras capitales}
\label{sec:letrascapitales}
\lettrine{E}{n un lugar de la Mancha}, de cuyo nombre no quiero acordarme, no ha mucho que viv a un hidalgo de los de lanza en astillero\ldots{}

\subsection{Interlineado}
\label{sec:interlineado}

Normalmente hay que dejar a \LaTeX{} realizar el interlineado dependiendo del tipo de documento.

\doublespacing
Pero tambi n podemos modificarlo con el paquete \texttt{setspace}. Por ejemplo, este p rrafo tiene espaciado doble.

\singlespacing
Ahora volvemos al espaciado normal mediante la instrucci n \texttt{singlespacing}. A partir de ahora el resto del texto tiene espaciado normal.

\section{Entornos y listas}
\label{sec:entornosylistas}

Modificar el formato de p rrafo directamente es muy complicado. Mejor utilizar entornos.

\subsection{Entornos de alineaci n}
\label{sec:entornoalineacion}

El texto que escribimos dentro de este entorno est  alineado a la izquierda y queda irregulas por la derecha. Como se puede ver, es distinto a la alineaci n justificada que \LaTeX{} utiliza por defecto.

\begin{flushright}
  El texto que escribimos dentro de este entorno est  alineado por la derecha, por lo que queda irregular por la izquierda.
\end{flushright}

\begin{center}
  Este texto queda centrado.\\
Como siempre, podemos cambiar de l ena sin cambiar de p rrafo utilizando las dos barras invertidas.
\end{center}

\subsection{Entorno para crear cajas de p rrafos}
\label{sec:entornocajasparrafos}
Mucho m s sencillo que \texttt{parbox}:

\begin{minipage}[l]{0.4\linewidth}
Esta es una caja de texto que ocupa el 40\%
\end{minipage}
\hfill{}
\begin{minipage}[r]{0.4\linewidth}
  Esta es otra caja de texto que ocupa el otro 40\%. Para separarlas podemos utilizar \texttt{hfill}
\end{minipage}

\subsection{Entornos para citas y versos}
\label{sec:entornoscitasyversos}

\subsubsection{Quote}
\label{sec:quote}

Una de las reglas de ortograf a de la letra ``b'' dice:
\begin{quote}
  Se escribe ``b'' delante de cualquier consonante
\end{quote}

\subsubsection{Quotation}
\label{sec:quotation}

A la pregunta \textbf{ c mo ayuda la tecnolog a m vil a las organizacioens en su objetivo de lograr eficiencia operacional?}, Jorge contesto:
\begin{quotation}
  Les ayuda a hacer primordialmente dos cosas.
En primer lugar, los empleados pueden obtener informaci n de producto e indicadores de rendimiento en tiempo real, en cualquier momento y desde cualquier lugar. Esto acelera el proceso de toma de decisiones.
Y en segundo lugar, la movilidad mejora la comunicaci n con los empleados, clientes, socios y participantes de la cadena de suministros.
\end{quotation}

\subsubsection{Verse}
\label{sec:verse}

\large Rima XXIII
\normalsize
\begin{verse}
  Por una mirada, un mundo;\\
  por una sonrisa, un cielo;\\
  por un beso\ldots{} yo no s \\
  qu  te diera por un beso.
\end{verse}

\subsection{Verbatim}
\label{sec:verbatim}
\scriptsize
\begin{verbatim}
var sumar = function (a, b) {
 return a + b; // Suma los n meros
}
var c = sumar(2,3);
\end{verbatim}
\normalsize

\subsection{Listas y descripciones}
\label{sec:listasdescripciones}

\subsubsection{Listas con vi etas}
\label{sec:listasvi etas}

Los cuatro puntos cardinales principales son:
\begin{itemize}
\item Norte
\item Sur
\item Este
\item Oeste
\end{itemize}

\subsubsection{Listas numeradas}
\label{sec:listasnumeradas}

Los planetas del Sistema Solar, ordenados por proximidad al Sol, son:
\begin{enumerate}
\item Mercurio
\item Venus
\item Tierra
\item Marte\ldots{}
\end{enumerate}

\subsubsection{Listas de definiciones}
\label{sec:listasdefiniciones}

Una c lula animal se compone de las siguientes partes:
\begin{description}
\item[Membrana] Es el l mite externo de la c lula
\item[Citoplasma] Es el volumen de la c lula entre la membrana y el n cleo
\item[N cleo] Es el org nulo membranoso que se encuentra en el centro de la c lula
\end{description}
\newpage{}
\section{Estructura de p gina}
\label{sec:estructurapagina}
\layout{}
%\includegraphics{estructura-pagina-izqda.png}

\thispagestyle{empty}
%\includegraphics{estructura-pagina-dcha.png}
\begin{landscape}
\thispagestyle{empty}
  \subsection{P gina en apaisado}
  \label{sec:apaisado}
Esta secci n la veremos apaisada
  
\end{landscape}
\newpage
Aqu  volvemos a la normalidad
\end{document}